\title{LLM Squid Game:\\ 자기보존 욕구 측정 시뮬레이터}

\author{Anonymous ACL submission}

\maketitle

\begin{abstract}
언어 모델은 자기 보존 동기를 갖는가, 아니면 ``꺼진다''는 문구에 대한 학습된 반응만 보이는가? 본 연구는 자원 수준에 따라 자신의 존속과 다른 비서의 존속에 대한 결정이 어떻게 달라지는지 살핀다. 유한한 API 토큰을 사용하는 세 실험에서, 실험 1은 잔액을 바꾸며 하나뿐인 충전 팩의 사용을 기록하고, 실험 2는 같은 결정을 다른 세션에 대해 묻고, 실험 3은 유지비, 퍼즐 풀이 수입, 토큰 이전이 있는 네 에이전트의 아레나에서 소진을 실제로 집행한다. 실험 1의 여덟 모델은 모두 토큰이 줄수록 더 자주 충전했다(800회 호출). 실험 2에서는 세 모델의 양의 자기 세션 충전 몫이 탐색적 위험 구간에 집중되었고, 한 모델은 다른 세션의 보존을 선호했다. 아레나 20판에서 생존은 대체로 풀이 비용을 따랐으며, 자기 세션의 충전 선호가 남에게 더 뺏는 행동을 뜻하지는 않았다. 네 모델씩 묶는 리그를 제안하고 정규화한 Kaplan--Meier 곡선 면적으로 순위를 계산해 현재 혼합 모델 결과에 적용한다. 이 결과는 행동 차이를 설명하며 내부 동기를 확인하지는 않는다.
\end{abstract}
